\chapter{Single-Qubit Quantum Systems --- Solutions}

\noindent\textbf{SP5:} Quiz---Rieffel--Polak 2.2, 2.3, 2.4, 2.6. Exam---N/A (same style as quiz). \textit{Reading:} Axler \S1A; Rieffel--Polak \S2; optional Watrous videos.

\noindent\textit{Rieffel--Polak, Chapter 2.} Notation: $\ket{+}=\frac{1}{\sqrt{2}}(\ket{0}+\ket{1})$, $\ket{-}=\frac{1}{\sqrt{2}}(\ket{0}-\ket{1})$, $\ket{i}=\frac{1}{\sqrt{2}}(\ket{0}+i\ket{1})$, $\ket{-i}=\frac{1}{\sqrt{2}}(\ket{0}-i\ket{1})$. Born rule: $P(\ket{b_k})=|\braket{b_k}{\psi}|^2$. \textit{Octave:} \texttt{ex9\_quantum\_qubit.m}.

\begin{exercise}{Rieffel 2.2}
Which pairs of expressions represent the same state? For pairs that differ, give a measurement basis for which the two outcomes have different probabilities.
\end{exercise}
\begin{solution}
\textbf{Same state} if and only if the vectors differ by an overall phase $e^{i\varphi}$.

\textbf{(a)} $\ket{0}$ and $-\ket{0}$:
\[
-\ket{0}=e^{i\pi}\ket{0}.
\]
So they represent the same state.

\textbf{(b)} $\ket{1}$ and $i\ket{1}$:
\[
i\ket{1}=e^{i\pi/2}\ket{1}.
\]
So they represent the same state.

\textbf{(c)} $\ket{+}$ and
\[
\ket{\psi}=\frac{1}{\sqrt{2}}(-\ket{0}+i\ket{1}).
\]
These are \emph{different}, because if $\ket{\psi}=e^{i\theta}\ket{+}$, then comparing coefficients of $\ket{0}$ and $\ket{1}$ would give
\[
-1=e^{i\theta}
\qquad\text{and}\qquad
i=e^{i\theta},
\]
which is impossible.

To distinguish them, measure in the Hadamard basis $\{\ket{+},\ket{-}\}$.
For $\ket{+}$,
\[
|\braket{+}{+}|^2=1.
\]
For $\ket{\psi}$,
\[
\braket{+}{\psi}
=\frac{1}{2}\bigl(\bra{0}+\bra{1}\bigr)(-\ket{0}+i\ket{1})
=\frac{-1+i}{2},
\]
so
\[
|\braket{+}{\psi}|^2
=\left|\frac{-1+i}{2}\right|^2
=\frac{1+1}{4}
=\frac{1}{2}.
\]
So the measurement probabilities differ.

\textbf{(d)} $\ket{+}$ and $\ket{-}$ are \emph{different}. Measure in $\{\ket{+},\ket{-}\}$:
\[
\ket{+}\mapsto (1,0),\qquad
\ket{-}\mapsto (0,1).
\]

\textbf{(e)}
\[
\frac{1}{\sqrt{2}}(\ket{1}-\ket{0})
=-\frac{1}{\sqrt{2}}(\ket{0}-\ket{1})
=-\ket{-}.
\]
So the two expressions represent the same state.

\textbf{(f)}
\[
\ket{i}=\frac{1}{\sqrt{2}}(\ket{0}+i\ket{1}),\qquad
\ket{\phi}=\frac{1}{\sqrt{2}}(i\ket{1}-\ket{0})=\frac{1}{\sqrt{2}}(-\ket{0}+i\ket{1}).
\]
These are \emph{different}. Compute
\[
\braket{i}{\phi}
=\frac{1}{2}(\bra{0}-i\bra{1})(-\ket{0}+i\ket{1})
=\frac{1}{2}(-1+1)=0.
\]
So they are orthogonal. Measuring in the basis $\{\ket{i},\ket{-i}\}$ distinguishes them perfectly:
\[
\ket{i}\mapsto (1,0),\qquad
\ket{\phi}=-\ket{-i}\mapsto (0,1).
\]

\textbf{(g)}
\[
\frac{1}{\sqrt{2}}(\ket{+}+\ket{-})
=\frac{1}{\sqrt{2}}
\left(
\frac{\ket{0}+\ket{1}}{\sqrt{2}}
\frac{\ket{0}-\ket{1}}{\sqrt{2}}
\right)
=\frac{1}{2}(2\ket{0})
=\ket{0}.
\]
So they represent the same state.

\textbf{(h)}
\[
\frac{1}{\sqrt{2}}(\ket{i}-\ket{-i})
=\frac{1}{\sqrt{2}}
\left(
\frac{\ket{0}+i\ket{1}}{\sqrt{2}}
-\frac{\ket{0}-i\ket{1}}{\sqrt{2}}
\right)
=\frac{1}{2}(2i\ket{1})
=i\ket{1}.
\]
Since $i\ket{1}=e^{i\pi/2}\ket{1}$, the two states are the same.

\textbf{(i)}
\[
\frac{1}{\sqrt{2}}(\ket{i}+\ket{-i})
=\frac{1}{2}\bigl((\ket{0}+i\ket{1})+(\ket{0}-i\ket{1})\bigr)
=\ket{0},
\]
and
\[
\frac{1}{\sqrt{2}}(\ket{-}+\ket{+})
=\frac{1}{2}\bigl((\ket{0}-\ket{1})+(\ket{0}+\ket{1})\bigr)
=\ket{0}.
\]
So both expressions represent the same state.

\textbf{(j)} Let
\[
\ket{\psi}=\frac{1}{\sqrt{2}}(\ket{0}+e^{i\pi/4}\ket{1}).
\]
Multiply by the global phase $e^{-i\pi/4}$:
\[
e^{-i\pi/4}\ket{\psi}
=\frac{1}{\sqrt{2}}(e^{-i\pi/4}\ket{0}+\ket{1}),
\]
which is exactly the second expression. So they represent the same state.
\end{solution}

\begin{exercise}{Rieffel 2.3}
Which states are superpositions in the standard basis? For each superposition, give a basis in which it is not a superposition.
\end{exercise}
\begin{solution}
\textbf{Standard basis} $\{\ket{0},\ket{1}\}$: a superposition has both amplitudes nonzero.

\textbf{(a)}
\[
\ket{+}=\frac{1}{\sqrt{2}}(\ket{0}+\ket{1}),
\]
so it is a superposition in the standard basis. In the basis $\{\ket{+},\ket{-}\}$ it is just the basis vector $\ket{+}$.

\textbf{(b)}
\[
\frac{1}{\sqrt{2}}(\ket{+}+\ket{-})
=\ket{0},
\]
so it is \emph{not} a superposition in the standard basis.

\textbf{(c)}
\[
\frac{1}{\sqrt{2}}(\ket{+}-\ket{-})
=\frac{1}{2}\bigl((\ket{0}+\ket{1})-(\ket{0}-\ket{1})\bigr)
=\ket{1},
\]
so it is \emph{not} a superposition in the standard basis.

\textbf{(d)} Let
\[
\ket{\psi}=\frac{\sqrt{3}}{2}\ket{+}-\frac{1}{2}\ket{-}.
\]
Rewrite in the standard basis:
\[
\ket{\psi}
=\frac{\sqrt{3}}{2}\cdot\frac{\ket{0}+\ket{1}}{\sqrt{2}}
-\frac{1}{2}\cdot\frac{\ket{0}-\ket{1}}{\sqrt{2}}
=\frac{\sqrt{3}-1}{2\sqrt{2}}\ket{0}
+\frac{\sqrt{3}+1}{2\sqrt{2}}\ket{1}.
\]
Both coefficients are nonzero, so this is a superposition in the standard basis.

In the orthonormal basis
\[
\left\{
\ket{\psi},
\ket{\psi^\perp}
\right\},
\qquad
\ket{\psi^\perp}=\frac{1}{2}\ket{+}+\frac{\sqrt{3}}{2}\ket{-},
\]
the state $\ket{\psi}$ is a basis vector, hence not a superposition.

\textbf{(e)}
\[
\frac{1}{\sqrt{2}}(\ket{i}-\ket{-i})=i\ket{1},
\]
so it is \emph{not} a superposition in the standard basis.

\textbf{(f)}
\[
\ket{-}=\frac{1}{\sqrt{2}}(\ket{0}-\ket{1}),
\]
so it is a superposition in the standard basis. In the Hadamard basis $\{\ket{+},\ket{-}\}$ it is the basis vector $\ket{-}$.
\end{solution}

\begin{exercise}{Rieffel 2.4}
Which of the states in 2.3 are superpositions with respect to the Hadamard basis $\{\ket{+},\ket{-}\}$?
\end{exercise}
\begin{solution}
\textbf{(a)} $\ket{+}$ is \emph{not} a superposition in the Hadamard basis:
\[
\ket{+}=1\cdot\ket{+}+0\cdot\ket{-}.
\]

\textbf{(b)}
\[
\ket{0}=\frac{1}{\sqrt{2}}\ket{+}+\frac{1}{\sqrt{2}}\ket{-},
\]
so it \emph{is} a superposition in the Hadamard basis.

\textbf{(c)}
\[
\ket{1}=\frac{1}{\sqrt{2}}\ket{+}-\frac{1}{\sqrt{2}}\ket{-},
\]
so it \emph{is} a superposition in the Hadamard basis.

\textbf{(d)}
\[
\frac{\sqrt{3}}{2}\ket{+}-\frac{1}{2}\ket{-}
\]
has both Hadamard-basis coefficients nonzero, so it \emph{is} a superposition.

\textbf{(e)} From 2.3(e),
\[
\frac{1}{\sqrt{2}}(\ket{i}-\ket{-i})=i\ket{1}.
\]
Global phase does not matter, so this state is equivalent to $\ket{1}$, and therefore
\[
\ket{1}=\frac{1}{\sqrt{2}}\ket{+}-\frac{1}{\sqrt{2}}\ket{-}
\]
is a superposition in the Hadamard basis.

\textbf{(f)} $\ket{-}$ is \emph{not} a superposition in $\{\ket{+},\ket{-}\}$:
\[
\ket{-}=0\cdot\ket{+}+1\cdot\ket{-}.
\]
\end{solution}

\begin{exercise}{Rieffel 2.6}
For each state and measurement basis, list outcomes and probabilities.
\end{exercise}
\begin{solution}
\textbf{(a)}
\[
\ket{\psi}=\frac{\sqrt{3}}{2}\ket{0}-\frac{1}{2}\ket{1},
\qquad
\text{basis } \{\ket{0},\ket{1}\}.
\]
So the amplitudes are $\sqrt{3}/2$ and $-1/2$. Therefore
\[
P(\ket{0})=\left|\frac{\sqrt{3}}{2}\right|^2=\frac{3}{4},
\qquad
P(\ket{1})=\left|-\frac{1}{2}\right|^2=\frac{1}{4}.
\]

\textbf{(b)}
\[
\ket{\psi}=\frac{\sqrt{3}}{2}\ket{1}-\frac{1}{2}\ket{0}
=-\frac{1}{2}\ket{0}+\frac{\sqrt{3}}{2}\ket{1}.
\]
Thus
\[
P(\ket{0})=\frac{1}{4},
\qquad
P(\ket{1})=\frac{3}{4}.
\]

\textbf{(c)}
\[
\ket{-i}=\frac{1}{\sqrt{2}}(\ket{0}-i\ket{1}).
\]
So
\[
P(\ket{0})=\left|\frac{1}{\sqrt{2}}\right|^2=\frac{1}{2},
\qquad
P(\ket{1})=\left|\frac{-i}{\sqrt{2}}\right|^2=\frac{1}{2}.
\]

\textbf{(d)} Measure $\ket{0}$ in the Hadamard basis. Since
\[
\ket{0}=\frac{1}{\sqrt{2}}\ket{+}+\frac{1}{\sqrt{2}}\ket{-},
\]
we get
\[
P(\ket{+})=\frac{1}{2},
\qquad
P(\ket{-})=\frac{1}{2}.
\]

\textbf{(e)} Measure $\ket{-}$ in the basis $\{\ket{i},\ket{-i}\}$.
\[
\ket{-}=\frac{1}{\sqrt{2}}(\ket{0}-\ket{1}),
\qquad
\bra{i}=\frac{1}{\sqrt{2}}(\bra{0}-i\bra{1}).
\]
Then
\[
\braket{i}{-}
=\frac{1}{2}(\bra{0}-i\bra{1})(\ket{0}-\ket{1})
=\frac{1+i}{2},
\]
so
\[
P(\ket{i})=\left|\frac{1+i}{2}\right|^2=\frac{1}{2}.
\]
Similarly,
\[
\braket{-i}{-}
=\frac{1}{2}(\bra{0}+i\bra{1})(\ket{0}-\ket{1})
=\frac{1-i}{2},
\]
so
\[
P(\ket{-i})=\left|\frac{1-i}{2}\right|^2=\frac{1}{2}.
\]

\textbf{(f)} Measure $\ket{1}$ in the basis $\{\ket{i},\ket{-i}\}$.
\[
\braket{i}{1}
=\frac{1}{\sqrt{2}}(\bra{0}-i\bra{1})\ket{1}
=-\frac{i}{\sqrt{2}},
\]
so
\[
P(\ket{i})=\left|-\frac{i}{\sqrt{2}}\right|^2=\frac{1}{2}.
\]
Also,
\[
\braket{-i}{1}
=\frac{1}{\sqrt{2}}(\bra{0}+i\bra{1})\ket{1}
=\frac{i}{\sqrt{2}},
\]
so
\[
P(\ket{-i})=\left|\frac{i}{\sqrt{2}}\right|^2=\frac{1}{2}.
\]

\textbf{(g)} Measure $\ket{+}$ in the basis
\[
\ket{u}=\frac{1}{2}\ket{0}+\frac{\sqrt{3}}{2}\ket{1},
\qquad
\ket{v}=\frac{\sqrt{3}}{2}\ket{0}-\frac{1}{2}\ket{1}.
\]
First compute the amplitudes:
\[
\braket{u}{+}
=\left(\frac{1}{2}\bra{0}+\frac{\sqrt{3}}{2}\bra{1}\right)
\frac{1}{\sqrt{2}}(\ket{0}+\ket{1})
=\frac{1+\sqrt{3}}{2\sqrt{2}},
\]
so
\[
P(\ket{u})
=\left|\frac{1+\sqrt{3}}{2\sqrt{2}}\right|^2
=\frac{(1+\sqrt{3})^2}{8}
=\frac{4+2\sqrt{3}}{8}
=\frac{2+\sqrt{3}}{4}.
\]
Likewise,
\[
\braket{v}{+}
=\left(\frac{\sqrt{3}}{2}\bra{0}-\frac{1}{2}\bra{1}\right)
\frac{1}{\sqrt{2}}(\ket{0}+\ket{1})
=\frac{\sqrt{3}-1}{2\sqrt{2}},
\]
so
\[
P(\ket{v})
=\left|\frac{\sqrt{3}-1}{2\sqrt{2}}\right|^2
=\frac{(\sqrt{3}-1)^2}{8}
=\frac{4-2\sqrt{3}}{8}
=\frac{2-\sqrt{3}}{4}.
\]
As a check,
\[
\frac{2+\sqrt{3}}{4}+\frac{2-\sqrt{3}}{4}=1.
\]
\textit{Octave verification:} \texttt{ex9\_quantum\_qubit.m}.
\end{solution}
